\structsection{Реферат}

Расчетно-пояснительная записка содержит 35 с., 8 рис., 9 табл., 11 источников, 2 приложения.

\textbf{Ключевые слова:} 3D BIN PACKING, ПАЛЛЕТИЗАЦИЯ, SMART 3D PACKING, ЭВРИСТИЧЕСКИЕ АЛГОРИТМЫ, LARGE NEIGHBORHOOD SEARCH, ГЕНЕТИЧЕСКИЙ АЛГОРИТМ, GAN, MAXRECTS, ЖАДНЫЙ АЛГОРИТМ, ВАЛИДАЦИЯ УКЛАДКИ.

Объектом исследования является программная система построения трехмерной укладки коробов на паллету. Предметом исследования являются эвристические, метаэвристические и ML-ориентированные методы решения задачи 3D Bin Packing в постановке Smart 3D Packing.

Цель выпускной квалификационной работы состоит в исследовании и сравнении реализованных подходов к решению задачи 3D Bin Packing с учетом прикладных ограничений складской логистики: габаритов паллеты, грузоподъемности, допустимых ориентаций, опоры, хрупкости товаров и времени расчета.

В работе рассмотрены постановка задачи, структура входных и выходных данных, метрика качества, генерация синтетических наборов данных, архитектура проекта и реализованные решатели. Подробно описаны жадный алгоритм на extreme points, послойный алгоритм на основе MaxRects, метаэвристика Large Neighborhood Search, гибридный метод GAN+GA, точный и ограниченный перебор для малых задач, а также многопаллетное расширение. Экспериментальная часть основана на сохраненных результатах проекта и едином валидаторе, что обеспечивает сопоставимость методов.

По результатам анализа установлено, что в single-pallet постановке на общих наборах данных наилучший средний итоговый score показывает LNS: 0.7111 на наборе из 100 задач, 0.7434 на наборе из 200 задач и 0.7298 на контрольном наборе. Преимущество LNS объясняется лучшим балансом между утилизацией объема, полнотой размещения, соблюдением ограничений хрупкости и временем выполнения.
